\scanpage{46}
\mspage{1}

I \quad $(N,M,d,\theta)=C$

\deleted{$\unsure{\theta:M\longrightarrow\operatorname{Aut}_{\mathrm{gr}}(N)}$}

\deleted{$\unsure{\beta:C'\longrightarrow C}$}

II \quad $(N',M',d',\theta')=C'_{\unsure{1}}$

III \quad $\text{\textcircled{$\beta$}}\left\{
\begin{array}{l}\pi_0(C)\simeq\pi_0(C'_{\unsure{1}})\\
\pi_1(C)\simeq\pi_1(C'_{\unsure{1}})
\end{array}\right.$ compatibles aux op. de $\pi_0$ sur $\pi_1$

\diagram{\begin{tikzcd}[ampersand replacement=\&,column sep=small,row sep=small]
\&N\cap N'=z=\pi_1\arrow[dl,dash]\arrow[dr,dash]\&\\
N\arrow[dr,dash]\arrow[ddr,dash,bend right=50,"M'"']\&\&N'\arrow[dl,dash]\arrow[ddl,dash,bend left=50,"M"]\\
\&NN'=H\arrow[d,dash]\&\\
\&G\&
\end{tikzcd}}

problème de reconstitution pour $K=G/z\simeq\operatorname{Im}(G\to M\times M')$
et $H=\operatorname{Im}(N\times N'\to G)$

\diagram{\begin{tikzcd}[ampersand replacement=\&,column sep=small]
\& (N\times N'\underset{\unsure{z}}{\longrightarrow}G,D_1,\Theta_1)\arrow[dl]\arrow[dr]\&\\
(N,M,d,\theta)\arrow[dr]\&\&(N',M',d',\theta')\arrow[dl]\\
\&(G,M\times M',D_2,\Theta_2)\&
\end{tikzcd}}
\[
(H,K,D,\Theta)
\]

\diagram{\begin{tikzcd}[ampersand replacement=\&,column sep=small]
1\arrow[r,leftrightarrow]\&z\arrow[r]\&H\arrow[r,dash,hook]\&G
\end{tikzcd}}
\figure{\begin{tikzpicture}[x=.8cm,y=.35cm]
\draw (1,1)--(1,1.5)--(3,1.5)--(3,1);\draw(0,0)--(0,-.5)--(2,-.5)--(2,0);
\end{tikzpicture}}

\diagram{\begin{tikzcd}[ampersand replacement=\&,column sep=large,row sep=large]
H\arrow[rrr,bend left=20]\&N\times N'\arrow[l]\arrow[r]\&M\times M'\arrow[d]\&K\arrow[l]\arrow[d]\\
\pi_1\arrow[u]\&\pi_1\times\pi_1\arrow[u]\arrow[l,"\text{épi}"']\&\pi_0\times\pi_0\arrow[r]\&\pi_0
\end{tikzcd}}
\quad cocart \qquad cart
\[
(\alpha,\alpha').(\lambda,\lambda')=(\alpha.\lambda,\alpha'.\lambda')
\]
\[
(\alpha,\alpha').(\lambda,-\lambda)=(\alpha.\lambda,-\alpha'\lambda)
\]

Obstruction dans $H^3(B_{\pi_0}/X,\pi_1)$

Indétermination dans $\begin{array}{c}H^2(B_{\pi_0}/X,\pi_1)\\ \Vert\\ \operatorname{Ext}_{\mathrm{gr}}(\pi_0,\pi_1)\end{array}$

\hrulefill

\deleted{\unsure{A}}
\[
\left\{\begin{array}{l}
N'\text{ central}\ \inserted{\text{i.e. }N'\subset\operatorname{Cent}(G)}
\Longleftrightarrow\uline{\theta':M'\longrightarrow\operatorname{Aut}_{\mathrm{gr}}(N')\text{ triviale}}\quad(\text{d'où }N'\text{ commut})\\
N\text{ cocentral}\ \inserted{\text{i.e. }N\supset DG}\Longleftrightarrow M'\text{ commutatif}
\end{array}\right.
\]

L'ensemble des deux conditions signifie que
la donnée de $C'$ se réduit à la donnée
de groupes comm. $N'$ et $M'$ et un
hom
\[
N'\xrightarrow{d'}M'
\]
sans plus. Cela implique que l'opération
de $\pi'_0$ sur $\pi'_1$ est triviale, donc itou pour
$\pi_0$ sur $\pi_1$ (condition à mettre sur $C$)
ceci dit, la donnée de
\[
\boxed{\begin{array}{cc}N'\xrightarrow{d'}M'&d'(x')\text{ \uline{central}}\\&N'\text{ commutatif}\end{array}}
\]
pour \textcircled{A} est sans condition
de compatibilité supplémentaire

\margin{\deleted{Lieber \unsure{Edgar},\par
Wie die lieben Leute sich miteinander vergnügen --- ob allein, zu zwein, zu drei, zu \unsure{vier|fünf} oder zu zwanzig}}
\margin{$\deleted{\unsure{H(\beta)}}$}
\mspage{45}
